\begin{theorem}[S-35: Prime-to-rank descent and the all-group rank restriction]
\label{thm:semireg-prime-to-rank-descent}
Let \((X,\theta)\) be a polarized complex abelian fourfold with a Rosati-invariant rational action of the real quadratic field \(F=\mathbf Q(\sqrt{\Delta_0})\), where \(\Delta_0>1\) is squarefree, of type \((2,2)\). Put \[
 D=\int_X\frac{\theta^4}{24},\qquad
 C=\frac{\theta_1^2+\theta_2^2}{2},\qquad
 \nu=\frac{\theta^4}{24},\qquad
 v_q=1-qC+q^2\nu.
 \tag{S35.1}
\] Let \(G\subset X\times\widehat X\) be a finite effective subgroup with a specified coherent action \[
 \Phi_g(K)=t_{a_g}^*K\otimes P_g,
\] and let a perfect object \(\mathcal E\) of positive virtual rank \(N\) have a coherent equivariant lift. Assume that for some positive rational \(q\) and rational degree-two class \(B\), \[
 e^B\operatorname{ch}(\mathcal E)=Nv_q.
 \tag{S35.2}
\] Let \(H\subseteq G\) be any subgroup whose order is coprime to \(N\).

\textbf{Descent assertion.} The projection \(H\to X\) is injective. There is an integer \(j\) such that \[
 \mathcal E\otimes\det(\mathcal E)^{-j}
           \simeq\pi^*\mathcal F
 \tag{S35.3}
\] for an ordinary perfect object \(\mathcal F\) of rank \(N\) on the abelian quotient \(\pi:X\to Y=X/H\). In particular, the scalar coherence obstruction on this subgroup can be removed; there is no residual Brauer twist on the asserted quotient object.

The character and integral-lattice consequence is \[
 t_H:=\frac{N^4Dq^2}{|H|}\in\mathbf Z_{>0},
 \qquad \Delta_0\mid t_H.
 \tag{S35.4}
\] This does not require \(\mathcal F\) to be simple or semiregular.

Suppose in addition that the given lift of \(\mathcal E\) is simple, gluable and fully semiregular in the sense used in S-22. Then \[
 k:=\frac{4DN^2q^2}{|G|}\in\{2,4\}.
\] The necessary rank conditions are \[
 \boxed{\Delta_0\mid N,\qquad k=2\Longrightarrow 2\mid N.}
 \tag{S35.5}
\] They hold for \textbf{every} effective finite coherent translation--tensor action. No integral maximal-order action, principal polarization or minimal Néron--Severi hypothesis is needed. Consequently rank one is impossible, and the rank restrictions previously obtained in S-30 for the nonequivariant case remain necessary for all such finite actions.
\end{theorem}
\begin{proof}
\paragraph{1. Determinants and injectivity of the translation projection}

Write \(L=\det(\mathcal E)\) and \(\ell=c_1(L)\). The degree-two part of (S35.2) forces \(B=-\ell/N\). For every \(g\in G\), the determinant of the equivariance isomorphism \(\Phi_g(\mathcal E)\simeq\mathcal E\) gives \[
 t_{a_g}^*L\otimes P_g^{\otimes N}\simeq L.
 \tag{S35.6}
\] This determinant formula uses the virtual rank and is valid for arbitrary perfect objects. The determinant of a scalar automorphism \(z\operatorname{id}_{\mathcal E}\) is \(z^N\).

If \(g\in H\) has \(a_g=0\), (S35.6) gives \(P_g^{\otimes N}\simeq\mathcal O_X\). The point \(P_g\in\widehat X\) also has order dividing \(|H|\), which is coprime to \(N\). Thus \(P_g=0\) in \(\widehat X\). Effectivity of the subgroup of \(X\times\widehat X\) then gives \(g=1\). Hence \(H\) is identified with a finite translation subgroup of \(X\); it acts freely, and \(Y=X/H\) is an abelian variety.

Let \(m\) be the exponent of \(H\). Since \(\gcd(N,m)=1\), choose \(j\in\mathbf Z\) such that \[
 Nj\equiv1\pmod{m^2},\qquad M:=1-Nj\in m^2\mathbf Z.
 \tag{S35.7}
\] For the trivial group, take \(m=1\); all statements below remain valid.

\paragraph{2. Removing the line-bundle factors}

Conjugate the given action by tensoring with \(L^{-j}\) and put \(\mathcal E'=\mathcal E\otimes L^{-j}\). The conjugated functor attached to \(g\in H\) is translation by \(a_g\) followed by tensoring with \[
 R_g=P_g\otimes(t_{a_g}^*L\otimes L^{-1})^{\otimes j}.
]
By (S35.6),
\[
 R_g\simeq P_g^{\otimes(1-Nj)}=P_g^{\otimes M}
                  \simeq\mathcal O_X,
 \tag{S35.8}
\] because the order of \(P_g\) divides \(m\) and \(m\mid M\). Choose trivializations of these line bundles. The functors are now the usual translations, but their multiplication isomorphisms may differ from the canonical ones by scalars. Their discrepancy is a two-cocycle \[
 \alpha\in Z^2(H,\mathbf C^*)
 \tag{S35.9}
\] with trivial coefficient action. The original coherence gives the cocycle identity. Changing the chosen trivializations changes \(\alpha\) by a coboundary. This is the usual \(H^2\) ambiguity of coherent actions with fixed underlying functors; natural degree-zero automorphisms of a translation Fourier--Mukai kernel are scalars on the connected projective variety \(X\).

The equivariance maps of \(\mathcal E'\) become an \(\alpha\)-projective translation linearization. We now show that the cohomology class of \(\alpha\) is zero.

\paragraph{3. The determinant kills the residual scalar coherence}

First, every line bundle \(L^{m^2}\) lies in the image of \[
 [m]^*:\operatorname{Pic}(X)\longrightarrow\operatorname{Pic}(X).
 \tag{S35.10}
\] Indeed, \([m]^*L\) and \(L^{m^2}\) have the same first Chern class, since \([m]^*\) acts as multiplication by \(m^2\) on \(H^2(X,\mathbf Z)\). Their ratio belongs to \(\operatorname{Pic}^0(X)\). Pullback \([m]^*\) is surjective on \(\operatorname{Pic}^0(X)\), being the dual isogeny to \([m]\). Thus one can adjust \(L\) by a degree-zero line bundle to obtain a line bundle \(Q\) with \([m]^*Q\simeq L^{m^2}\). This argument does not assert \([m]^*L\simeq L^{m^2}\) for a nonsymmetric \(L\).

Because \(M\) is divisible by \(m^2\), the line bundle \[
 \det(\mathcal E')=L^{\otimes M}
\] is also pulled back by \([m]\). The translation subgroup \(H\) is contained in \(X[m]\), so \([m]\) factors through \(X/H\). Consequently this determinant line bundle has an honest \(H\)-linearization.

On the other hand, taking the determinant of the \(\alpha\)-projective maps on \(\mathcal E'\) gives a projective linearization of the same line bundle with scalar cocycle \(\alpha^N\). Compare it with an honest linearization. Their ratios are global line-bundle automorphisms, hence scalars, so the comparison says \[
 N[\alpha]=0\quad\text{in }H^2(H,\mathbf C^*).
 \tag{S35.11}
\] The standard restriction--corestriction identity says that this positive-degree cohomology group is killed by \(|H|\). Since \(\gcd(N,|H|)=1\), equation (S35.11) forces \([\alpha]=0\). Choosing a coboundary trivialization converts the action and the equivariance maps into an ordinary coherent translation action and an honest \(H\)-linearization of \(\mathcal E'\).

For the precise scalar convention, if \(\alpha(g,h)\alpha(gh,k)=\alpha(h,k)\alpha(g,hk)\), then \(b(g)=\prod_{h\in H}\alpha(g,h)\) satisfies \(\alpha(g,h)^{|H|}=b(g)b(h)b(gh)^{-1}\). This is the degree-two instance of the cited annihilation identity; it also directly checks the use of the trivial coefficient action here.

\paragraph{4. Ordinary descent and its normalized character}

Let \(\pi:X\to Y=X/H\) be the quotient isogeny. Finite étale descent gives a perfect object \(\mathcal F\) with \(\pi^*\mathcal F\simeq\mathcal E'\). One can realize it concretely as \[
 \mathcal F=(\pi_*\mathcal E')^H.
 \tag{S35.12}
\] The pushforward along a finite étale map preserves perfectness, and the averaging idempotent splits a perfect direct summand. Finite flat base change gives \(\pi^*\pi_*\mathcal E'\simeq\bigoplus_{h\in H}t_h^*\mathcal E'\); the honest linearization identifies the invariant summand with \(\mathcal E'\), proving the assertion. In particular \(\operatorname{rk}\mathcal F=N\).

There is a rational ample class \(\theta_Y\) with \(\pi^*\theta_Y=\theta\). It can be obtained by dividing the norm polarization by \(|H|\). The rational real multiplication descends by the isogeny and remains Rosati-invariant of type \((2,2)\). Let \(C_Y,\nu_Y\) be its corresponding classes. Then \[
 \pi^*C_Y=C,\quad \pi^*\nu_Y=\nu,\quad
 \int_Y\nu_Y=\frac D{|H|}.
 \tag{S35.13}
\] The first Chern class of \(\mathcal E'\) is \(M\ell\). Twisting by a line bundle leaves the determinant-normalized character unchanged, so \[
 e^{-c_1(\mathcal E')/N}\operatorname{ch}(\mathcal E')=Nv_q.
\] Since \(\pi^*\) is injective on rational cohomology, the descended object satisfies \[
 e^{-c_1(\mathcal F)/N}\operatorname{ch}(\mathcal F)
       =N(1-qC_Y+q^2\nu_Y).
 \tag{S35.14}
\] No semiregularity assertion about this partial quotient object is needed.

Choose a positive integer \(c\) such that \(\theta'_Y=c\theta_Y\) is an integral polarization. Put \(q'=q/c^2\). Then \(C'_Y=c^2C_Y\), \(\nu'_Y=c^4\nu_Y\), so the right side of (S35.14) is exactly \(Nv_{q'}\) for the integral polarization \(\theta'_Y\). Its volume is \(D'_Y=c^4D/|H|\).

For every positive-rank perfect object with this normalized character, the second-Chern identity used in S-30 gives the integral class \[
 N^2q'C'_Y
     =N c_2(\mathcal F)
        -(N-1)\frac{c_1(\mathcal F)^2}{2}.
\] The term \(c_1(\mathcal F)^2/2\) is integral on a torus. Applying the integral four-form obstruction S-30 to this class yields \[
 N^4D'_Y(q')^2\in\mathbf Z_{>0},
 \qquad \Delta_0\mid N^4D'_Y(q')^2.
\] The rescaling cancels exactly: {[} N\^{}4D'\_Y(q')\^{}2=\frac{N^4Dq^2}{|H|}=t\_H. {]} This proves (S35.4).

\paragraph{5. Removing all primes not dividing the rank}

The finite group \(G\subset X\times\widehat X\) is abelian. Let \(H\) be the product of its primary subgroups at primes not dividing \(N\), and write \[
 G_N=G/H.
]
Every prime divisor of $|G_N|$ divides $N$, while $\gcd(|H|,N)=1$, so the preceding argument applies. The semiregular numerical identity of S-22 converts (S35.4) to
\[
 t_H=\frac{N^2 k |G_N|}{4}\in\Delta_0\mathbf Z_{>0}.
 \tag{S35.15}
\] If an odd prime \(p\mid\Delta_0\) did not divide \(N\), then it would divide none of \(N,k,|G_N|,4\); equation (S35.15) would be a \(p\)-adic unit, contradicting its divisibility by \(\Delta_0\).

Suppose \(N\) is odd. Then \(|G_N|\) is odd. If \(k=2\), the right side of (S35.15) is a half-integer, contradicting integrality. If \(k=4\) and \(2\mid\Delta_0\), it is an odd integer, contradicting divisibility by \(\Delta_0\). Every prime factor of the squarefree integer \(\Delta_0\) therefore divides \(N\), and the \(k=2\) case always requires even \(N\). This is (S35.5).
\end{proof}
